\begin{theorem}[Theorem~\ref{thm:random_poly} restated]
For every~$d\in \N$ and~$\sigma,\eps\in (0,1]$, there exist constants~$\kappa = \kappa(d, \sigma)>0$ and~$R = R(d,\sigma,\eps)\in\N$ such that the following holds.
For every natural rank function~$\rank: \Pol_{\leq d}(\F^{\infty}, \F^k) \to \R$ and every map~$\phi \in \Pol_{\leq d}(\F^n, \F^k)$ with~$\rank(\phi) \geq R$, we have
\beqn
\Pr_{I\sim [n]_{\sigma}}\big[ \rank(\phi_{|I}) \geq \kappa\cdot \rank(\phi)\big] \geq 1 - \eps.
\eeqn
\end{theorem}

The proof of this theorem proceeds differently from the proof of Theorem~\ref{thm:random_tensor}, which gives the analogous result for tensors.
The reason for this is that the natural choice of function $f:\bset{n}\to \R_+$ to which one might apply the concentration inequality in Lemma~\ref{lem:subadditive}, given by $f(\one_I) = \rank(\phi_{|I})$ (for some rank function~$\rank$), might fail to be sub-additive as a function on the boolean hypercube while~$\rank$ is sub-additive as a natural rank function.
For instance, consider (in dimension~$k=1$) a $2n$-variate polynomial $p(x_1,\dots,x_n,y_1,\dots,y_n)$ such that each monomial contains both an~$x$ variable and a~$y$ variable.
Then the rank of~$p$ restricted only to its~$x$ variables is zero, and similarly for the restriction to the~$y$ variables, while~$\rank(p)$ can be of order~$\Theta(n)$.
This shows that~$\rank(p_{|I \cup J})$ can be much larger than $\rank(p_{|I}) + \rank(p_{|J})$, and the argument used in the tensor case breaks down.


\subsection{The proof in expectation}

The first step in our proof of Theorem~\ref{thm:random_poly} is to show that the desired result is true \emph{in expectation}, rather than with high probability.
More precisely, we wish to show an inequality of the form
\beqn
\Exp_{I\sim [n]_{\sigma}} \rank(\phi_{|I}) \geq c(\sigma, d) \rank(\phi)
\eeqn
which is valid for all degree-$d$ maps $\phi\in \Pol_{\leq d}(\F^n, \F^k)$ and all natural rank functions.

Towards proving the inequality above by induction on the degree, it will be convenient to generalize it to polynomial maps over \emph{commutative rings} rather than only fields.
The reason for doing so is that, by considering some subset of variables $\{x_a: a\in A\}$ to be constants, we might be able to decrease the degree of the associated polynomial map over the remaining variables $\{x_b: b\notin A\}$.

More formally, given a polynomial map $\phi\in \Pol_{\leq d}(\F^n, \F^k)$ and two disjoint subsets $A, B \subset [n]$, denote by $\phi[A, B]$ the sum of the monomials in~$\phi$ involving at least one variable from each~$A$ and~$B$.
Note that $\phi[A, B]$ can be regarded as a polynomial map in three different ways:
either as having coefficients in~$\F$ and variables in $\{x_i: i\in [n]\}$;
or having coefficients in the ring $\F[x_b: b\in B]$ and variables in $\{x_a: a\in A\}$;
or having coefficients in $\F[x_a: a\in A]$ and variables in $\{x_b: b\in B\}$.
The advantage of these last two representations is that the degree of $\phi[A, B]$ \emph{decreases}, since every monomial will contain formal variables which are now elements of the ring (and are thus regarded as constants).

For a commutative ring~$R$, denote by $R[x_1, \dots, x_n]_{\leq d}$ the set of formal polynomials in variables $x_1, \dots, x_n$ with coefficients over~$R$ and degree at most~$d$.
For a bipartition $[n] = A\uplus B$, define the polynomial ring $R_A = R[x_a\st a\in A]$ and consider the natural identification of $R[x_1,\dots,x_n]$ and $R_A[x_b\st b\in B]$.
Under this identification, a function $\rank: (R[x_1, \dots, x_n])^k \to \R_+$ induces a function $\rank:(R_A[x_b\st b\in B])^k\to \R_+$ in a way that preserves symmetry, sub-additivity and monotonicity under restrictions.


\begin{lemma}[Linear expectation] \label{lem:expectation}
For every~$\sigma\in (0, 1]$ and every integer~$d \geq 1$ there exists a constant~$c(\sigma, d) > 0$ such that the following holds.
Let~$n,k$ be positive integers,~$R$ be a commutative ring and $\rank: (R[x_1, \dots, x_n]_{\leq d})^k \to \R_+$ be symmetric, sub-additive and monotone under restrictions.
Then, for any map $\phi\in (R[x_1, \dots, x_n]_{\leq d})^k$ we have
\beqn
\Exp_{J\sim [n]_{\sigma}} \rank(\phi_{|J}) \geq c(\sigma, d) \rank(\phi).
\eeqn
\end{lemma}

\begin{proof}
For simplicity, we will prove the result in the special case where $\sigma=1/2$.
Since $\phi_{|I\cap J} = (\phi_{|I})_{|J}$ and $I\cap J \sim [n]_{1/2^{j+1}}$ when $I\sim [n]_{1/2}$, $J\sim [n]_{1/2^j}$, the general case easily follows from this special case by taking constant
\beqn
c(\sigma, d) = c(1/2, d)^{\lceil\log(1/\sigma)\rceil}
\eeqn
and using monotonicity under restrictions.

Denote $c(d) := c(1/2, d)$ for ease of notation.
The proof will proceed by induction on the degree~$d$, with the base case where $d=0$ holding trivially with $c(0) = 1$.
For the inductive step, let
\beqn
c(d) = \frac{c(d-1)^2}{20}
\eeqn
and consider two different possibilities.

First, suppose that for every bipartition $[n] = A\uplus B$ we have that
\beqn
\max \big\{\rank(\phi_{|A}),\, \rank(\phi_{|B})\big\} \geq 2c(d) \rank(\phi).
\eeqn
Using a similar coupling argument as was used in the proof of Lemma~\ref{lem:subadditive}, we get that
\begin{align*} \Pr_{I\sim[n]_{1/2}}\big[\rank\big(\phi_{|I}\big) \geq 2c(d)\rank(\phi)\big] \geq \frac{1}{2},
\end{align*}
which implies
\begin{equation*}
    \Exp_{I\sim [n]_{1/2}}\rank\big(\phi_{|I}\big) \geq c(d)\rank(\phi)
\end{equation*}
as desired.

The second possibility is that there exists a bipartition $[n] = A\uplus B$ such that
\beq \label{eq:second_case}
\rank(\phi_{|A}) < 2c(d) \rank(\phi) \quad \text{and} \quad \rank(\phi_{|B}) < 2c(d) \rank(\phi).
\eeq
Fix such a bipartition for the rest of the proof.
For any given subsets $I\subseteq A$, $J\subseteq B$, let $\phi[I, J]$ be the sum of the monomials in~$\phi$ involving at least one variable from each~$I$ and~$J$.
It is easy to see that
\beqn
\phi[I, J] = \phi_{|I\cup J} - \phi_{|I} - \phi_{|J} + \phi(0).
\eeqn
Define the commutative rings $R_A := R[x_a: a\in A]$ and $R_B := R[x_b: b\in B]$, and note that $\phi[I, J]$ can be seen both as an element of $(R_A[x_b: b\in B])^k$ and as an element of $(R_B[x_a: a\in A])^k$.
Moreover, the degree of $\phi[I, J]$ is at most $d-1$ in each of these representations.

From the identity $\phi = \phi_{|A} + \phi_{|B} - \phi(0) + \phi[A, B]$ and sub-additivity, it follows that
\beqn
\rank(\phi[A, B]) \geq \rank(\phi) - \rank(\phi_{|A}) - \rank(\phi_{|B}) - \rank(-\phi(0)).
\eeqn
Since $\rank(-\phi(0)) = \rank(\phi(0)) = \rank(\phi_{|\emptyset})$, we conclude from monotonicity under restrictions and our assumption~\eqref{eq:second_case} that
\beq\label{eq:phiAB}
\rank(\phi[A, B]) > \rank(\phi) - 6 c(d) \rank(\phi) > \rank(\phi)/2.
\eeq
In an analogous way, we obtain the inequalities
\begin{align} \label{eq:phi_IJ}
    \rank(\phi_{|I\cup J}) &\geq \rank(\phi[I, J]) - \rank(\phi_{|I}) - \rank(\phi_{|J}) - \rank(\phi(0)) \\
    &> \rank(\phi[I, J]) - 6 c(d) \rank(\phi) \notag
\end{align}
valid for all $I\subseteq A$, $J\subseteq B$.

Applying the inductive hypothesis to $\phi[A, B] \in (R_B[x_a: a\in A]_{\leq d-1})^k$ we obtain
\beqn
\Exp_{I\sim A_{1/2}} \rank(\phi[I, B]) = \Exp_{I\sim A_{1/2}} \rank(\phi[A, B]_{|I}) \geq c(d-1) \rank(\phi[A, B]).
\eeqn
Moreover, for any fixed $I\subseteq A$, the inductive hypothesis applied to the map $\phi[I, B] \in (R_A[x_b: b\in B]_{\leq d-1})^k$ gives
\beqn
\Exp_{J\sim B_{1/2}} \rank(\phi[I, J]) = \Exp_{J\sim B_{1/2}} \rank(\phi[I, B]_{|J}) \geq c(d-1) \rank(\phi[I, B]).
\eeqn
Combining the two inequalities above we conclude that
\beqn
\Exp_{I\sim A_{1/2},\, J\sim B_{1/2}} \rank(\phi[I, J]) \geq c(d-1)^2 \rank(\phi[A, B]),
\eeqn
and thus by~\eqref{eq:phiAB} we have
$\Exp_{I\sim A_{1/2},\, J\sim B_{1/2}} \rank(\phi[I, J]) > c(d-1)^2 \rank(\phi)/2$.
Together with equation~$\eqref{eq:phi_IJ}$, we conclude that
\begin{align*}
    \Exp_{I\sim A_{1/2},\, J\sim B_{1/2}}\rank(\phi_{|I\cup J})
    &> \Exp_{I\sim A_{1/2},\, J\sim B_{1/2}} \rank(\phi[I, J]) - 6c(d) \rank(\phi) \\
    &> c(d-1)^2 \rank(\phi)/2 - 6c(d) \rank(\phi) \\
    &> c(d) \rank(\phi),
\end{align*}
which is precisely what we wanted to prove.
\end{proof}




\subsection{Monotone functions on the hypercube and boosting}

Our next result is a lemma which allows us to boost the probability of some events (such as having high rank under random restrictions) from~$\eps$ to~$1-\eps$ by paying a relatively small price.

It will again be convenient to take a more abstract approach and deal with Boolean functions rather than restrictions of polynomials.
For a Boolean function~$g: \{0,1\}^n \to \{0,1\}$, the total influence of~$g$ under distribution~$\pi_{\sigma}^n$ is given by
\beqn
\mathbf{I}^{(\sigma)}(g)
=
\sum_{i=1}^n \Exp_{x\sim\pi_\sigma^n}\big|g(x) - g(x^i)\big|,
\eeqn
where~$x^i$ differs from~$x$ only in the~$i$th coordinate.
Denote by~$\mu_{\sigma}[g] := \Exp_{x\sim \pi_{\sigma}^n} \big[g(x)\big]$ its expectation.

\begin{lemma}[Boosting lemma] \label{lem:boost}
For every~$\sigma>0$ there is a constant~$C_{\sigma}>0$ such that the following holds.
Let~$n\in \N$, $f: \{0,1\}^n \to \R_+$ be a monotone~$1$-Lipschitz function and~$\eps\in (0, 1/2]$.
If~$r\geq 1$ satisfies
\beqn
\Pr_{x\sim \pi_{\sigma}^n} \big[f(x) \geq r\big] > \eps,
\eeqn
then~$\Pr_{x\sim \pi_{1.01\sigma}^n} \big[f(x) \geq r - C_{\sigma}^{1/\eps^2}\big] > 1-\eps$.
\end{lemma}

\begin{proof}
Consider the Boolean function~$g(x) := 1 \big[f(x) \geq r\big]$.
Note that there exists~$q \in [\sigma,\, 1.01\sigma]$ such that~$\mathbf{I}^{(q)}(g) \leq 100/\sigma$, as otherwise by the Margulis-Russo formula we would have
\beqn
\mu_{1.01\sigma}[g] = \mu_{\sigma}[g] + \int_{\sigma}^{1.01\sigma} \frac{d\mu_q[g]}{dq} dq \geq \int_{\sigma}^{1.01\sigma} \mathbf{I}^{(q)}(g) dq > 1.
\eeqn

Let~$\gamma \in (0, \eps/2)$ be a constant to be chosen later.
By Friedgut's Junta Theorem~\cite{Friedgut1998}, there is a~$C(q)^{100/\gamma \sigma}$-junta~$h: \{0,1\}^n \to \{0,1\}$ such that
\beqn
\Pr_{x\sim \pi_q^n} \big[g(x) \neq h(x)\big] < \gamma.
\eeqn
Let~$J \subseteq [n]$, $|J| \leq C(q)^{100/\gamma \sigma}$, be the set of variables on which~$h$ depends.
Then
\begin{align*}
    \mu_q[h] \geq \mu_q[g] - \gamma \geq \mu_{\sigma}[g] - \gamma > \eps/2,
\end{align*}
which implies~$\Pr_{z\sim \pi_q^J} \big[h(z) = 1\big] > \eps/2$.
Since
\beqn
\Exp_{z\sim \pi_q^J} \Pr_{y\sim \pi_q^{J^c}} \big[g(y,z) \neq h(z)\big] < \gamma,
\eeqn
it follows that
\beqn
\Pr_{z\sim \pi_q^J} \big[\Pr_{y\sim \pi_q^{J^c}} \big[g(y,z) \neq h(z)\big] \geq \eps\big] < \gamma/\eps.
\eeqn
Taking~$\gamma = \eps^2/2$, we conclude there exists~$z\in \{0,1\}^J$ such that~$h(z)=1$ and
\beqn
    \Pr_{y\sim \pi_q^{J^c}} \big[g(y,z) = 0\big] = \Pr_{y\sim \pi_q^{J^c}} \big[g(y,z) \neq h(z)\big] < \eps.
\eeqn

Since~$f$ is~$1$-Lipschitz by assumption, we have
\beqn
    f(y,z) \geq r \implies f(y,z') \geq r - |J| \quad \forall z'\in \{0,1\}^J,
\eeqn
and thus by monotonicity
\begin{align*}
    \Pr_{x\sim \pi_{1.01\sigma}^n} \big[f(x) \geq r - |J|\big]
    &\geq \Pr_{x\sim \pi_q^n} \big[f(x) \geq r - |J|\big] \\
    &= \Exp_{z'\sim \pi_q^J} \Pr_{y\sim \pi_q^{J^c}} \big[f(y,z') \geq r - |J|\big] \\
    &\geq \max_{z\in \{0,1\}^J} \Pr_{y\sim \pi_q^{J^c}} \big[f(y,z) \geq r\big] \\
    &> 1-\eps.
\end{align*}
This is precisely what we wanted to prove, with constant
$$C_{\sigma} = \sup_{q\in [\sigma,\, 1.01\sigma]} C(q)^{200/\sigma}.$$
Note that the above supremum is finite:
in \cite[Section~10.3]{ODonnell:2014} it is shown that we can take
\beqn
C(q) = \max\big\{ q^{-c},\, (1-q)^{-c}\big\}
\eeqn
for some absolute constant $c>0$.
\end{proof}


\subsection{The Random Restriction Theorem}

Our main result, Theorem~\ref{thm:random_poly}, follows easily from the lemmas given above.
Recall that we wish to show that
\beqn
\Pr_{J\sim [n]_{\sigma}}\big[ \rank(\phi_{|J}) \geq \kappa(d, \sigma)\cdot \rank(\phi)\big] \geq 1 - \eps
\eeqn
whenever~$\rank(\phi) \geq R(d, \sigma, \eps)$, for some well-chosen constants~$R(d, \sigma, \eps)$ and $\kappa(d, \sigma) > 0$.

\begin{proof}[ of Theorem~\ref{thm:random_poly}]
Applying Lemma \ref{lem:expectation} with~$\sigma$ substituted by~$0.9\sigma$, we obtain
\beqn
\Exp_{J\sim [n]_{0.9\sigma}} \rank(\phi_{|J}) \geq c(0.9\sigma, d) \rank(\phi).
\eeqn
Since~$\rank(\phi_{|J}) \leq \rank(\phi)$ for all subsets~$J \subseteq [n]$, we conclude that
\beqn
\Pr_{J\sim [n]_{0.9\sigma}} \bigg[ \rank(\phi_{|J}) \geq \frac{c(0.9\sigma, d)}{2} \rank(\phi) \bigg] \geq \frac{c(0.9\sigma, d)}{2}.
\eeqn

By possibly decreasing~$\eps$ a little, we may assume that~$\eps < c(0.9\sigma, d)/2$.
Denote
\beqn
\kappa(d, \sigma) = \frac{c(0.9\sigma, d)}{4} \quad \text{and} \quad R(d, \sigma, \eps) = \frac{4 C_{0.9\sigma}^{1/\eps^2}}{c(0.9\sigma, d)},
\eeqn
where~$C_{0.9\sigma} > 0$ is the constant guaranteed by the boosting lemma, Lemma \ref{lem:boost}.
Applying the boosting lemma to the function~$J \mapsto \rank(\phi_{|J})$ with~$r = c(0.9\sigma, d) \rank(\phi)/2$, we get
\begin{align*}
    \Pr_{J\sim [n]_{0.9\sigma}} \big[ \rank(\phi_{|J}) \geq r \big] &\geq c(0.9\sigma, d)/2 > \eps \\
    &\implies \Pr_{J\sim [n]_{\sigma}} \big[ \rank(\phi_{|J}) \geq r - C_{0.9\sigma}^{1/\eps^2} \big] > 1-\eps.
\end{align*}
Since~$c(0.9\sigma, d) \rank(\phi)/2 - C_{0.9\sigma}^{1/\eps^2} \geq \kappa(\sigma, d) \rank(\phi)$ whenever $\rank(\phi) \geq R(d, \sigma, \eps)$, the result follows.
\end{proof}
